\section{Predicciones}
Registradas el 2026-11-07 (16:42, hora de Lima) por correo institucional, antes de cualquier
simulación de los incisos E1.1--E1.6. Los resultados se toman de \texttt{results.json} (\texttt{valores\_informe}).

\begin{center}\small
\begin{tabularx}{\linewidth}{@{}>{\raggedright\arraybackslash}p{1.1cm} >{\raggedright\arraybackslash}p{1.9cm} >{\raggedright\arraybackslash}X >{\raggedright\arraybackslash}p{3.3cm} >{\raggedright\arraybackslash}p{3.6cm}@{}}
\toprule
\texttt{id} & Registro (H1) & Enunciado (tipo, base) & Resultado & Acierto \\
\midrule
E1-P1 & 2026-11-07 &
El certificado elipsoidal del diseño base no garantiza $|\alpha_0|>0.055$~rad (resto en cero); el Monte
Carlo no lineal recupera desde ángulos al menos 5 veces mayores ($\ge0.27$~rad). (LIMITE, TEORIA) &
Certificado $\V{e16_corte_alpha}$~rad; recuperación desde $\V{e16_al_emp}$~rad (estado real) y $\V{e16_rad_oficial_exito}$~rad con éxito en el simulador oficial. &
SI: el certificado queda bajo la franja $V_{\max}/|K_2|=\V{e16_franja}$~rad y la región real es $\V{e16_P1_razon_of}$ a $\V{e16_P1_razon_id}$ veces mayor. \\
E1-P2 & 2026-11-07 &
Cumplir C2 ($|\theta|\le\pi/2$, $\le5\,\%$ saturado) exige $R\ge10R_0$, o bien $|K_2|\le40$~V/rad.
(ORDEN, TEORIA) &
El diseño base ($R_0$, $K_2=\V{e15_P2_K2_base}$~V/rad) ya cumple C2: saturado $\V{esc_C2_sat}\,\%$, $\max|\theta|=\V{esc_C2_th}$~rad. Con $Q_0$ y $R\ge\V{e15_P2_m_falla}R_0$ ($K_2\le\V{e15_P2_K2_falla}$) el brazo supera $\pi/2$. &
NO: la tendencia es la opuesta. La saturación dura poco ($\V{esc_C2_sat}\,\%$ del tiempo) y un regulador más suave deja que el brazo se aleje más. \\
E1-P3 & 2026-11-07 &
Al reducir $R$ de 1 a $10^{-8}$, el error de la Hamiltoniana crece más rápido que el de Schur; con
$R=10^{-8}$: Hamiltoniana $>10^{-6}$, Schur $<10^{-9}$. (SIGNO, TEORIA + IA) &
Hamiltoniana $\V{e13_P3_ham}$, Schur $\V{e13_P3_schur}$; pendientes $\V{e13_P3_pend_ham}$ y $\V{e13_P3_pend_schur}$. &
NO: la tendencia fue la opuesta. \texttt{eig} balancea $H$ por defecto y \texttt{schur} no; con el
mismo balanceo ambos métodos dan el mismo error (E1.3). \\
\bottomrule
\end{tabularx}
\end{center}
